\section*{Chapter \ref{ch:b3:hilbert} --- Hilbert Spaces}

\begin{solution}{exo:b3:hilbert:1}
(a) Real: expand $\norm{x \pm y}^2 = \norm x^2 \pm 2\langle
x,y\rangle + \norm y^2$ and subtract. Complex (\omterm{def:b3:hilbert:inner}{inner product}
linear in the second slot): expanding as above,
\[
\langle x, y\rangle = \frac14\sum_{k=0}^{3}
\iu^k\,\bigl\|\iu^kx + y\bigr\|^2,
\]
each term contributing $\iu^k\cdot2\operatorname{Re}\bigl(
(-\iu)^k\langle x,y\rangle\bigr)$, whose sum is $4\langle
x,y\rangle$ (check the four values of $k$; the $\sum\iu^k
(\norm x^2 + \norm y^2) = 0$).

(b) $L^1$: $f = \mathbf 1_{\intcc0{1/2}}$, $g = \mathbf
1_{\intcc{1/2}1}$: $\norm{f\pm g}_1^2 = 1$ each, sum $2$;
$2\norm f_1^2 + 2\norm g_1^2 = 1 \neq 2$. Sup norm: $f =
\mathbf 1$, $g(t) = t$ on $\intcc01$: $\norm{f + g}_\infty^2 +
\norm{f-g}_\infty^2 = 4 + 1 = 5 \neq 4 = 2 + 2$. Failing the
parallelogram law, these norms are induced by no \omterm{def:b3:hilbert:inner}{inner product}
(which would force it by direct expansion).
\end{solution}

\begin{solution}{exo:b3:hilbert:2}
(a) $p(f) = \bigl(\int_0^1f\bigr)\mathbf 1$: indeed $f - \int f
\perp$ constants ($\int(f - \int f)c = 0$). The best constant
approximation of $f$ in mean square is its \emph{average} ---
the first instance of conditional expectation
(\cref{ch:b3:probability}).

(b) $p(f) = f\,\mathbf 1_{\intcc{1/2}1}$: the difference
$f\mathbf 1_{\intcc0{1/2}}$ is orthogonal to every $g$
vanishing on $\intcc0{1/2}$.

(c) $d^2 = \bigl\|x - \tfrac12\bigr\|_2^2 = \int_0^1(x -
\tfrac12)^2\dd x = \tfrac1{12}$: $d = \frac1{2\sqrt3}$.
\end{solution}

\begin{solution}{exo:b3:hilbert:3}
(a) The equivalence is \cref{thm:b3:hilbert:decomposition}
($\bar F = (F^\perp)^\perp$, and $\bar F = H \iff F^\perp =
\{0\}$). Example: the space $F$ of finite sequences is dense
in $\ell^2$ (truncation) and proper: $F^\perp = \{0\}$ yet $F
\neq \ell^2$ --- for a non-closed subspace, $H = F \oplus
F^\perp$ fails blatantly ($F \oplus \{0\} \neq H$).

(b) $\abs{\langle x_n, y_n\rangle - \langle x, y\rangle} \leq
\abs{\langle x_n - x, y_n\rangle} + \abs{\langle x, y_n -
y\rangle} \leq \norm{x_n - x}\sup_n\norm{y_n} + \norm
x\,\norm{y_n - y} \to 0$ (convergent sequences are bounded).
Used: in \cref{thm:b3:hilbert:parseval}(1) to see $x - \lim
S_N \perp e_k$, and in \cref{thm:b3:hilbert:decomposition}
to see $F^\perp = \bar F^{\,\perp}$.
\end{solution}

\begin{solution}{exo:b3:hilbert:4}
$e_0 = \frac1{\sqrt2}$. Next, $x \perp \mathbf 1$ already
($\int_{-1}^1x = 0$), and $\int_{-1}^1x^2 = \frac23$: $e_1 =
\sqrt{\tfrac32}\,x$. Then $x^2 - \langle e_0, x^2\rangle e_0 =
x^2 - \frac13$ (and $\perp e_1$ by parity), with
\[
\int_{-1}^1\Bigl(x^2 - \frac13\Bigr)^2\dd x = \frac25 -
\frac49 + \frac29 = \frac{8}{45}:
\qquad e_2 = \sqrt{\tfrac{45}8}\,\Bigl(x^2 - \frac13\Bigr).
\]
Comparison: $P_0 = 1$, $P_1 = x$, $P_2 = \frac{3x^2 - 1}2$, and
$\sqrt{n + \tfrac12}\,P_n$ gives $\frac1{\sqrt2}$,
$\sqrt{\frac32}x$, $\sqrt{\frac52}\,\frac{3x^2-1}2 =
\sqrt{\frac{45}8}\bigl(x^2 - \frac13\bigr)$: exactly $e_0, e_1,
e_2$.
\end{solution}

\begin{solution}{exo:b3:hilbert:5}
For $f(t) = t$: $c_0 = 0$ and, integrating by parts, $c_n =
\frac{\iu(-1)^n}{n}$ for $n \neq 0$: $\abs{c_n}^2 =
\frac1{n^2}$. \omterm{thm:b3:hilbert:parseval}{Parseval}:
\[
\frac1{2\pi}\int_{-\pi}^\pi t^2\dd t = \frac{\pi^2}3
= \sum_{n\neq0}\frac1{n^2} = 2\sum_{n\geq1}\frac1{n^2}
\ \Longrightarrow\ \sum_{n\geq1}\frac1{n^2} = \frac{\pi^2}6 .
\]
For $f(t) = t^2$: $c_0 = \frac{\pi^2}3$, $c_n =
\frac{2(-1)^n}{n^2}$ ($n \ne 0$). \omterm{thm:b3:hilbert:parseval}{Parseval}:
\[
\frac1{2\pi}\int_{-\pi}^{\pi}t^4\dd t = \frac{\pi^4}5
= \frac{\pi^4}9 + \sum_{n\neq0}\frac4{n^4}
\ \Longrightarrow\
\sum_{n\geq1}\frac1{n^4} = \frac18\Bigl(\frac{\pi^4}5 -
\frac{\pi^4}9\Bigr) = \frac{\pi^4}{90} .
\]
No piecewise-$\mathcal C^1$ caveats are needed:
\cref{thm:b3:hilbert:fourier} covers every $L^2$ function.
\end{solution}

\begin{solution}{exo:b3:hilbert:6}
(a) $\varphi(f) = \int_0^{1/2}f = \langle\mathbf
1_{\intcc0{1/2}},\ f\rangle$: the representative is $a =
\mathbf 1_{\intcc0{1/2}}$, and $\norm\varphi = \norm a_2 =
\frac1{\sqrt2}$ (\cref{thm:b3:hilbert:riesz}).

(b) Take the tent functions $f_n$ with peak $1$ at
$\frac12$ and support of width $\frac2n$: $f_n(\tfrac12) = 1$
while $\norm{f_n}_2^2 \leq \frac2n \to 0$: no constant $C$
can give $\abs{f(\frac12)} \leq C\norm f_2$. Point evaluation
is meaningless in $L^2$ --- elements are classes modulo null
sets --- and this computation is the quantitative reason.
\end{solution}

\begin{solution}{exo:b3:hilbert:7}
(a) Let $M = \sup_k\norm{x_k}$. The scalar sequences
$(\langle e_n, x_k\rangle)_k$ are bounded by $M$: a diagonal
extraction yields $x_{k_j}$ with $\langle e_n, x_{k_j}\rangle
\to \gamma_n$ for every $n$. For each $N$: $\sum_{n\leq
N}\abs{\gamma_n}^2 = \lim_j\sum_{n\leq N}\abs{\langle e_n,
x_{k_j}\rangle}^2 \leq M^2$ (Bessel), so $(\gamma_n) \in
\ell^2$ and $x = \sum_n\gamma_ne_n \in H$
(\cref{thm:b3:hilbert:parseval}(3)). For $y \in H$:
\[
\abs{\langle y, x_{k_j} - x\rangle}
\leq \Bigl|\sum_{n\leq N}\overline{c_n(y)}\bigl(\langle e_n,
x_{k_j}\rangle - \gamma_n\bigr)\Bigr|
+ 2M\Bigl(\sum_{n>N}\abs{c_n(y)}^2\Bigr)^{1/2},
\]
using the expansion $\langle y, z\rangle =
\sum\overline{c_n(y)}c_n(z)$ and Cauchy--Schwarz on the tail;
choose $N$ then $j$: weak convergence to $x$.

(b) $\langle y, e_n\rangle = c_n(y) \to 0$ for every $y$
($\ell^2$ tails): $e_n \rightharpoonup 0$, yet $\norm{e_n} =
1$: the norm is not weakly \omterm{def:b3:topology:continuity}{continuous}. In (a): $\norm x^2 =
\sum\abs{\gamma_n}^2 \leq \liminf_j\norm{x_{k_j}}^2$ (finite
sections and Bessel again): weak limits can only lose norm.
\end{solution}

\begin{solution}{exo:b3:hilbert:8}
For fixed $y$, $x \mapsto \langle y, Tx\rangle$ is a
\omterm{def:b3:topology:continuity}{continuous} linear functional; Riesz gives a unique $T^*y$ with
$\langle y, Tx\rangle = \langle T^*y, x\rangle$ for all $x$
--- conjugating, $\langle Tx, y\rangle = \langle x,
T^*y\rangle$. Uniqueness makes $T^*$ linear;
\[
\norm{T^*y} = \sup_{\norm x = 1}\abs{\langle T^*y, x\rangle}
= \sup_{\norm x=1}\abs{\langle y, Tx\rangle}
\leq \vertiii T\,\norm y,
\]
so $\vertiii{T^*} \leq \vertiii T$, and $T^{**} = T$ gives
equality. Shift: $\langle Sx, y\rangle = \sum_{n\geq1}
x_n\bar y_{n+1} = \langle x, S^*y\rangle$ with $(S^*y)_n =
y_{n+1}$: the backward shift. Kernel--image: $T^*y = 0$ iff
$\langle x, T^*y\rangle = 0$ for all $x$ iff $\langle Tx,
y\rangle = 0$ for all $x$ iff $y \perp \operatorname{im}T$:
$\ker T^* = (\operatorname{im}T)^\perp$; taking $\perp$ and
using \cref{thm:b3:hilbert:decomposition},
$\overline{\operatorname{im}T} = (\ker T^*)^\perp$.
\end{solution}

\begin{solution}{exo:b3:hilbert:9}
If $a(u, \cdot) = \varphi$: for any $w$,
\[
J(u + w) - J(u) = a(u, w) - \varphi(w) + \tfrac12a(w,w)
= \tfrac12a(w,w) \geq \tfrac\alpha2\norm w^2,
\]
strictly positive for $w \neq 0$: $u$ is the unique minimizer.
Conversely, at a minimizer the function $t \mapsto J(u + tw)$
(a quadratic polynomial in $t$) has vanishing derivative at
$0$: $a(u, w) - \varphi(w) = 0$ for every $w$. Projection
re-derived: for a closed subspace $F$, apply Lax--Milgram on
the \omterm{def:b3:hilbert:inner}{Hilbert space} $F$ with $a(u,v) = \langle u, v\rangle$
($M = \alpha = 1$) and $\varphi(v) = \langle x, v\rangle$: a
unique $p \in F$ with $\langle p, v\rangle = \langle x,
v\rangle$ for all $v \in F$, i.e.\ $x - p \perp F$ --- and by
the symmetric case, $p$ minimizes $\frac12\norm v^2 - \langle
x, v\rangle = \frac12\norm{v - x}^2 - \frac12\norm x^2$ over
$F$: the projection.
\end{solution}

\begin{solution}{exo:b3:hilbert:10}
Normalization: $\int h_n^2 = 2^j\cdot 2^{-j} = 1$.
Orthogonality: two distinct Haar functions have either
disjoint (\omterm{def:b3:topology:interior}{interiors} of) supports (product zero a.e.), or the
support of the finer is contained in a half-interval where the
coarser is constant --- then the integral of the product is
that constant times $\int h_{\text{finer}} = 0$; against $h_0
= \mathbf 1$, again $\int h_n = 0$. Totality: the span of
$\{h_0, \dots, h_{2^J-1}\}$ consists of step functions on the
dyadic grid of step $2^{-J}$; both spaces have dimension
$2^J$ and the Haar functions are independent (orthonormal):
the span is \emph{all} such step functions. Dyadic step
functions are dense in $L^2(\intcc01)$: \omterm{def:b3:lebesgue:simple}{simple functions} are
dense (\cref{thm:b3:lp:density}(1)), \omterm{def:b3:lebesgue:measurable}{measurable} sets are
approximated by finite unions of intervals
(\cref{exo:b3:measure:7}), and intervals by dyadic ones
(endpoints move by $\leq 2^{-J}$). By
\cref{thm:b3:hilbert:parseval}, the Haar system is a \omterm{def:b3:hilbert:onb}{Hilbert
basis}.
\end{solution}

\begin{solution}{exo:b3:hilbert:11}
(i) $\Rightarrow$ (ii): for the orthogonal projection,
$\langle Px, y\rangle = \langle Px, Py\rangle = \langle x,
Py\rangle$ (insert the decompositions $x = Px + (x - Px)$
etc.\ and kill cross terms). (ii) $\Rightarrow$ (iii):
$\norm{Px}^2 = \langle P^2x, x\rangle = \langle Px, x\rangle
\leq \norm{Px}\norm x$, so $\vertiii P \leq 1$, and $Pu = u$
on the nonzero image: $= 1$. (iii) $\Rightarrow$ (i): $H =
\operatorname{im}P \oplus \ker P$ (\omterm{def:b3:galois:algebraic}{algebraically}, from $P^2
= P$); take $u = Pu \in \operatorname{im}P$, $v \in \ker P$,
$t \in \R$: $\norm{P(u + tv)}^2 = \norm u^2$ must be $\leq
\norm{u + tv}^2 = \norm u^2 + 2t\operatorname{Re}\langle u,
v\rangle + t^2\norm v^2$ for every $t$, forcing
$\operatorname{Re}\langle u, v\rangle = 0$ (compare the
linear terms as $t \to 0^\pm$); replacing $v$ by $\iu v$
kills the imaginary part too: $\operatorname{im}P \perp
\ker P$, which is exactly orthogonality of the projection.
Example: $P(x, y) = (x + y, 0)$ on $\R^2$: $P^2 = P$,
image the $x$-axis, kernel the line $y = -x$, and
$\vertiii P = \sup\frac{\abs{x+y}}{\norm{(x,y)}} = \sqrt2$
(attained at $(1,1)/\sqrt2$): an oblique projection has norm
$> 1$. (For the record, (ii) also gives (i) directly:
$\ker P = \ker P^* = (\operatorname{im}P)^\perp$ by
\cref{exo:b3:hilbert:8}.)
\end{solution}

\begin{solution}{exo:b3:hilbert:12}
(a) For unitary $U$: $\norm{Ux - x}^2 = 2\norm x^2 -
2\operatorname{Re}\langle Ux, x\rangle$ and $\norm{U^*x -
x}^2 = 2\norm x^2 - 2\operatorname{Re}\langle x, Ux\rangle$:
the two vanish together, so $\ker(U - I) = \ker(U^* - I)$.
Then, using $\ker T^* = (\operatorname{im}T)^\perp$
(\cref{exo:b3:hilbert:8}) with $T = U - I$ and $T^* = U^* -
I$:
\[
\overline{\operatorname{im}(U - I)} = \bigl(\ker(U^* -
I)\bigr)^\perp = F^\perp .
\]

(b) On $F$: $U^kx = x$, so $A_nx = x$. For $x = (U - I)y$:
$A_nx = \frac1n(U^ny - y)$, of norm $\leq \frac2n\norm y \to
0$. For $x$ in the \omterm{def:b3:topology:interior}{closure}: given $\varepsilon$, pick $x' =
(U - I)y$ with $\norm{x - x'} < \varepsilon$; since
$\vertiii{A_n} \leq \frac1n\sum\vertiii{U^k} = 1$,
$\norm{A_nx} \leq \norm{A_n(x - x')} + \norm{A_nx'} \leq
\varepsilon + o(1)$.

(c) Decompose $x = Px + (x - Px)$ with $Px \in F$ and $x -
Px \in F^\perp = \overline{\operatorname{im}(U - I)}$ (part
(a)): $A_nx = Px + A_n(x - Px) \to Px + 0$.

(d) In the Fourier basis $e_m(x) = \eu^{2\iu\pi mx}$: $Ue_m
= \eu^{2\iu\pi m\alpha}e_m$, so $Ue_m = e_m$ iff $m\alpha
\in \Z$ iff $m = 0$ ($\alpha$ irrational): $F = \C\mathbf 1$
and $Pf = \langle\mathbf 1, f\rangle\mathbf 1 = \int_0^1f$.
The theorem reads $\frac1n\sum_{k<n}f(\cdot + k\alpha) \to
\int_0^1f$ in $L^2(\R/\Z)$: the \omterm{def:b3:groups:action}{orbit} averages of an
irrational rotation equidistribute --- the $L^2$ shadow of
Weyl's equidistribution theorem, obtained by pure Hilbert
geometry.
\end{solution}

\begin{solution}{pb:b3:hilbert:1}
\textbf{1.} Gram--Schmidt guarantees
$\operatorname{Vect}(p_0, \dots, p_n) =
\operatorname{Vect}(1, \dots, t^n) = \R_n[t]$ and $p_n \perp
p_k$ ($k < n$), hence $p_n \perp \R_{n-1}[t]$. The $p_k$, of
strictly increasing degrees, are independent: a basis.

\textbf{2.} $t\,p_n$ is monic of degree $n + 1$: expand
$t\,p_n = p_{n+1} + \sum_{k\leq n}c_kp_k$ with $c_k =
\langle p_k, tp_n\rangle/\norm{p_k}^2$. For $k \leq n - 2$:
$\langle p_k, tp_n\rangle = \langle tp_k, p_n\rangle = 0$
(degree $k + 1 < n$). So $tp_n = p_{n+1} + a_np_n +
b_np_{n-1}$, the stated recurrence, with
\[
b_n = \frac{\langle p_{n-1}, tp_n\rangle}{\norm{p_{n-1}}^2}
= \frac{\langle tp_{n-1}, p_n\rangle}{\norm{p_{n-1}}^2}
= \frac{\langle p_n + (\text{lower}),\
p_n\rangle}{\norm{p_{n-1}}^2}
= \frac{\norm{p_n}^2}{\norm{p_{n-1}}^2} > 0 .
\]

\textbf{3.} Let $t_1 < \dots < t_m$ be the points \omterm{def:b3:topology:interior}{interior} to
$I$ where $p_n$ changes sign, and $q = \prod_{i\leq m}(t -
t_i)$ (with $q = 1$ if $m = 0$). Then $p_nq$ has constant sign
on $I$ and is not a.e.\ zero: $\int_Ip_nq\,w \neq 0$. If $m <
n$, this contradicts $p_n \perp \R_{n-1}[t]$. So $m = n$:
$p_n$ has $n$ distinct \omterm{def:b3:topology:interior}{interior} roots (it has at most $n$
roots in all).

\textbf{4.} $(t^2 - 1)^n$ has degree $2n$; $n$ derivatives
leave degree $n$, with leading coefficient
$\frac{(2n)(2n-1)\cdots(n+1)}{2^nn!} =
\frac{(2n)!}{2^n(n!)^2}$. For $\deg Q < n$, integrate by parts
$n$ times: all boundary terms contain a derivative of order
$< n$ of $(t^2-1)^n$, which vanishes at $\pm1$ (root of order
$n$); after $n$ steps the integrand carries $Q^{(n)} = 0$.

\textbf{5.} With $u = (t^2 - 1)^n$:
\[
(2^nn!)^2\norm{P_n}^2 = \int_{-1}^1(u^{(n)})^2
= (-1)^n\int_{-1}^1 u\,u^{(2n)}
= (2n)!\int_{-1}^1(1 - t^2)^n\dd t ,
\]
($u^{(2n)} = (2n)!$; boundary terms vanish as in question 4).
And $\int_{-1}^1(1-t^2)^n\dd t = B(\tfrac12, n+1) =
\frac{\Gamma(\frac12)\Gamma(n+1)}{\Gamma(n + \frac32)} =
\frac{2\cdot4^n(n!)^2}{(2n+1)!}$
(\cref{exo:b3:product:8}). Combining: $\norm{P_n}^2 =
\frac{2}{2n + 1}$.

\textbf{6.} Polynomials are $\norm\cdot_\infty$-dense in
$\mathcal C(\intcc{-1}1)$ (Weierstrass,
\cref{cor:b3:complete:weierstrass}), \omterm{def:b3:topology:continuity}{continuous} functions are
$L^2$-dense (\cref{thm:b3:lp:density}), and
$\norm\cdot_2 \leq \sqrt2\norm\cdot_\infty$: polynomial spans
are total, so the normalized $P_n$ form a \omterm{def:b3:hilbert:onb}{Hilbert basis}.
Expansion of $\abs t$: the coefficient against $P_0$ is
$\frac{\langle P_0, \abs t\rangle}{\norm{P_0}^2} = \frac12$;
against $P_1$: $0$ (parity); against $P_2$:
$\frac{\int_{-1}^1\abs t\,\frac{3t^2-1}2\dd t}{2/5} =
\frac{1/4}{2/5} = \frac58$. Best quadratic approximation:
\[
\abs t \approx \frac12 + \frac58\,P_2(t) = \frac{3}{16} +
\frac{15}{16}\,t^2 .
\]

\textbf{7.} From $\frac{\dd^{n+1}}{\dd t^{n+1}}\eu^{-t^2} =
\frac{\dd^n}{\dd t^n}(-2t\,\eu^{-t^2})$ and Leibniz, $H_{n+1}
= 2tH_n - H_n'$; induction gives degree $n$ and leading
coefficient $2^n$. For $m < n$, integrate by parts $n$ times
in $\int H_m H_n\eu^{-t^2} = (-1)^n\int H_m\,\bigl(\eu^{-t^2}
\bigr)^{(n)}$: boundary terms (polynomial $\times$
$\eu^{-t^2}$) vanish at $\pm\infty$, leaving $\int
H_m^{(n)}\,\eu^{-t^2} = 0$. For $m = n$: $H_n^{(n)} = 2^nn!$,
so $\norm{H_n}_w^2 = 2^nn!\int\eu^{-t^2} = 2^nn!\sqrt\pi$.

\textbf{8.} Let $f \in L^2(\R, \eu^{-t^2}\dd t)$ be orthogonal
to every polynomial, and $g = f\eu^{-t^2}$. Then $g \in
L^1$: $\int\abs f\eu^{-t^2} \leq \bigl(\int\abs
f^2\eu^{-t^2}\bigr)^{1/2}\bigl(\int\eu^{-t^2}\bigr)^{1/2}$
(Cauchy--Schwarz). For $\xi \in \R$, expand $\eu^{-\iu\xi t}$:
the partial sums are dominated since
\[
\sum_k\frac{\abs\xi^k}{k!}\int\abs f\,\abs t^k\eu^{-t^2}\dd t
\leq \Bigl(\int \abs f^2\eu^{-t^2}\Bigr)^{1/2}
\sum_k\frac{\abs\xi^k}{k!}\Bigl(\int
t^{2k}\eu^{-t^2}\Bigr)^{1/2} < \infty
\]
(the last series converges: $\int t^{2k}\eu^{-t^2} =
\Gamma(k+\frac12) \leq k!\,\sqrt\pi$, so the terms are
$O(\abs\xi^k/\sqrt{k!})$). Term-by-term integration
(\cref{cor:b3:lebesgue:additivity} applied to the absolute
series, then Fubini for series) gives
\[
\int_\R g(t)\,\eu^{-\iu\xi t}\dd t
= \sum_k\frac{(-\iu\xi)^k}{k!}\int f(t)\,t^k\,\eu^{-t^2}\dd t
= 0 ,
\]
each integral being $\langle t^k, f\rangle_w$-type $= 0$. By
the admitted injectivity of the Fourier transform
(\cref{ch:b3:fouriertransform}), $g = 0$ a.e., so $f = 0$
a.e.: the Hermite family (whose spans are the polynomials) is
total.

\textbf{9.} Exactness to degree $n - 1$: for such $P$, $P =
\sum_iP(t_i)\ell_i$ exactly, so $\int Pw = \sum_iP(t_i)\int
\ell_iw = Q(P)$. Degree $\leq 2n - 1$: divide $P = qp_n + r$,
$\deg q \leq n - 1$, $\deg r \leq n-1$; then $\int Pw = \int
qp_nw + \int rw = 0 + Q(r)$ ($p_n \perp \R_{n-1}[t]$), while
$Q(P) = \sum_iw_i\bigl(q(t_i)\,p_n(t_i) + r(t_i)\bigr) = Q(r)$
since the nodes are the roots of $p_n$. Equal.

\textbf{10.} $\ell_i^2$ has degree $2n - 2 \leq 2n - 1$ and
$\ell_i^2(t_j) = \delta_{ij}$: $0 < \int\ell_i^2w = Q(\ell_i^2)
= w_i$. Polya (\cref{exo:b3:banach:9}, transported to $I$ with
weight): condition (i) holds --- each polynomial is integrated
exactly once $2n - 1 \geq$ its degree; condition (ii):
$\sum_i\abs{w_{i}} = \sum_iw_i = Q(\mathbf 1) = \int_Iw$,
bounded: $Q_n(f) \to \int fw$ for every $f \in \mathcal C(I)$,
$I$ \omterm{def:b3:topology:compact}{compact}.

\textbf{11.} Monic $p_2 = t^2 - \frac13$ (from
\cref{exo:b3:hilbert:4}): nodes $\pm\frac1{\sqrt3}$. Weights:
$\ell_1(t) = \frac{t - \frac1{\sqrt3}}{-\frac2{\sqrt3}}$, and
$w_1 = \int_{-1}^1\ell_1 = 1$; by symmetry $w_2 = 1$.
Exactness: $\int 1 = 2 = 1 + 1$; $\int t = 0 =
-\frac1{\sqrt3} + \frac1{\sqrt3}$; $\int t^2 = \frac23 =
\frac13 + \frac13$; $\int t^3 = 0$. The two-point trapezoid
rule (nodes $\pm1$, weights $1, 1$) is exact only to degree
$1$: on $t^2$ it returns $2$ instead of $\frac23$. Same cost,
two extra degrees of exactness: the payoff of orthogonal
nodes.

\textbf{12.} From $\cos(n{+}1)\theta + \cos(n{-}1)\theta =
2\cos\theta\cos n\theta$: $T_{n+1} = 2tT_n - T_{n-1}$ with
$T_0 = 1$, $T_1 = t$; induction gives polynomials of degree
$n$ with leading coefficient $2^{n-1}$ ($n \geq 1$).
Substituting $t = \cos\theta$ ($w(t)\dd t \mapsto
\dd\theta$): $\langle T_m, T_n\rangle_w =
\int_0^\pi\cos m\theta\cos n\theta\,\dd\theta = 0$ for $m
\neq n$, $= \pi$ for $m = n = 0$, $= \frac\pi2$ otherwise
(product-to-sum). Degrees and pairwise orthogonality
identify the $T_n$ with the Gram--Schmidt output up to
scalars; a Chebyshev expansion of $f$ is exactly the
Fourier cosine series of $\theta \mapsto f(\cos\theta)$.

\textbf{13.} $T_n(t) = 0$ iff $\cos n\theta = 0$ iff
$\theta = \frac{(2k-1)\pi}{2n}$: the $n$ distinct roots
$t_k = \cos\frac{(2k-1)\pi}{2n} \in \intoo{-1}1$. Extrema:
$\abs{T_n} \leq 1$ on $\intcc{-1}1$, with $T_n(s_j) =
(-1)^j$ at the $n + 1$ points $s_j = \cos\frac{j\pi}n$:
\omterm{prop:b3:galois:perfect}{perfect} equioscillation.

\textbf{14.} $2^{1-n}T_n$ is monic with sup-norm $2^{1-n}$.
If a monic $P$ of degree $n$ had $\sup\abs P < 2^{1-n}$, the
difference $D = 2^{1-n}T_n - P$ would have degree $\leq n -
1$ (leading terms cancel) yet alternate in sign at $s_0 >
\dots > s_n$ (there $2^{1-n}T_n = \pm2^{1-n}$ dominates
$P$): at least $n$ zeros --- $D \equiv 0$, contradiction.
For uniqueness at equality, the same $D$ satisfies
$(-1)^jD(s_j) \geq 0$; a nonzero polynomial of degree $\leq
n-1$ cannot have $n$ weakly alternating extremal
constraints without $n$ roots counted properly (if
$D(s_j) = 0$ for some \omterm{def:b3:topology:interior}{interior} $s_j$, that zero is double
in the counting since $D$ keeps a sign locally): again $D
\equiv 0$.

\textbf{15.} The Lagrange error formula (Rolle, Year~2)
gives $f - L_nf = \frac{f^{(n)}(\xi_t)}{n!}\,\omega(t)$, so
the uniform error is at most $\frac{\norm{f^{(n)}}_\infty}
{n!}\,\sup\abs\omega$, and $\omega$ is monic of degree $n$:
by question 14, $\sup_{\intcc{-1}1}\abs\omega \geq 2^{1-n}$
with equality iff the nodes are the Chebyshev roots. Hence
the optimal bound $\norm{f - L_nf}_\infty \leq
\frac{\norm{f^{(n)}}_\infty}{2^{n-1}n!}$. With equally
spaced nodes, $\sup\abs\omega$ is exponentially larger near
the endpoints, and interpolating even $\frac1{1 + 25t^2}$
diverges there as $n \to \infty$ (Runge's phenomenon);
Chebyshev nodes are the cure.

\textbf{16.} Differentiating $T_n(\cos\theta) = \cos
n\theta$: $T_n'(\cos\theta) = \frac{n\sin n\theta}
{\sin\theta}$, which tends to $n^2$ as $\theta \to 0$ and to
$(-1)^{n+1}n^2$ as $\theta \to \pi$: $\abs{T_n'(\pm1)} =
n^2$. At \omterm{def:b3:topology:interior}{interior} points, $\abs{T_n'(t)} \leq
\frac{n}{\sqrt{1 - t^2}} = O(n)$: the quadratic blow-up
lives only at the edges (Bernstein's \omterm{def:b3:topology:interior}{interior} bound versus
Markov's global one).

\textbf{17.} Let $\theta_k = \frac{(2k-1)\pi}{2n}$ and $S_j
= \sum_{k=1}^n\cos(j\theta_k)$ for $1 \leq j \leq n-1$. Then
\[
S_j = \operatorname{Re}\Bigl[\eu^{\iu j\pi/2n}
\sum_{k=0}^{n-1}\eu^{\iu jk\pi/n}\Bigr]
= \operatorname{Re}\Bigl[\eu^{\iu j\pi/2n}\,
\frac{\eu^{\iu j\pi} - 1}{\eu^{\iu j\pi/n} - 1}\Bigr] .
\]
For $j$ even the numerator vanishes: $S_j = 0$. For $j$ odd
the numerator is $-2$, and $\eu^{\iu j\pi/n} - 1 =
\eu^{\iu j\pi/2n}\cdot2\iu\sin\frac{j\pi}{2n}$, so the whole
expression is $\frac{-2}{2\iu\sin(j\pi/2n)} =
\frac{\iu}{\sin(j\pi/2n)}$: purely imaginary, $S_j = 0$
again. Hence the equal-weight rule $\frac\pi n\sum_kf(t_k)$
integrates $T_0$ ($\sum w_i = \pi = \int w$) and kills $T_1,
\dots, T_{n-1}$ exactly as $\int T_jw = 0$ does: it is exact
to degree $n - 1$. Weights exact to degree $n-1$ at given
nodes are unique (Lagrange basis): the Gauss weights are all
$\frac\pi n$. For $n = 3$: nodes $\pm\frac{\sqrt3}2, 0$ and
\[
\int_{-1}^1\frac{f(t)}{\sqrt{1 - t^2}}\,\dd t \approx
\frac\pi3\Bigl[f\Bigl(\tfrac{\sqrt3}2\Bigr) + f(0) +
f\Bigl(-\tfrac{\sqrt3}2\Bigr)\Bigr],
\]
exact through degree $5$.

\textbf{18.} For monic $P$ of degree $n$: $P = p_n + r$
with $r \in \R_{n-1}[t]$, and $p_n \perp \R_{n-1}[t]$
(question 1), so $\norm P^2 = \norm{p_n}^2 + \norm r^2 \geq
\norm{p_n}^2$, with equality iff $r = 0$: $p_n$ is the
orthogonal projection residue of $t^n$ onto
$\R_{n-1}[t]^\perp$, i.e.\ the monic polynomial closest to
the subspace it must avoid. Chebyshev's $2^{1-n}T_n$ answers
the same question for the sup-norm: least deviation from
zero, once in $L^2(w)$, once in $L^\infty$.

\textbf{19.} Write $K_n(x, y) =
\sum_{k=0}^n\frac{p_k(x)p_k(y)}{h_k}$. Base $n = 0$: $(x -
y)\frac1{h_0} = \frac{p_1(x)\cdot1 - 1\cdot p_1(y)}{h_0}$
since $p_1 = t - a_0$. Step: assuming the identity for $n -
1$,
\[
(x - y)\,K_n(x,y) = \frac{p_n(x)p_{n-1}(y) -
p_{n-1}(x)p_n(y)}{h_{n-1}} +
\frac{(x - y)\,p_n(x)p_n(y)}{h_n} ;
\]
substitute $x\,p_n(x) = p_{n+1}(x) + a_np_n(x) +
b_np_{n-1}(x)$ and $y\,p_n(y) = p_{n+1}(y) + a_np_n(y) +
b_np_{n-1}(y)$ in the second term: the $a_n$ contributions
cancel, and the $b_n = \frac{h_n}{h_{n-1}}$ contributions
cancel the induction term; what survives is
$\frac{p_{n+1}(x)p_n(y) - p_n(x)p_{n+1}(y)}{h_n}$. The
confluent form follows by letting $y \to x$ (both sides are
polynomials in $y$).

\textbf{20.} The confluent form gives $p_{n+1}'p_n -
p_n'p_{n+1} = h_n\sum_{k\leq n}\frac{p_k^2}{h_k} \geq
\frac{h_n}{h_0} > 0$ everywhere. At a root $x_0$ of
$p_{n+1}$: $p_{n+1}'(x_0)\,p_n(x_0) > 0$, so $p_n(x_0) \neq
0$ (no common roots). Between consecutive roots $x_0 < x_1$
of $p_{n+1}$ (all simple, Part I), $p_{n+1}'$ has opposite
signs, hence so does $p_n$: a root of $p_n$ lies in each of
the $n$ gaps --- and that exhausts its $n$ roots:
interlacing.

\textbf{21.} Expanding $D_n(t) = \det(tI_n - J_n)$ along the
last row: $D_n = (t - a_{n-1})D_{n-1} - b_{n-1}D_{n-2}$,
with $D_0 = 1$, $D_1 = t - a_0$: the recurrence and seeds of
the monic $p_n$, so $D_n = p_n$. Roots of $p_n$ =
eigenvalues of the symmetric $J_n$: real, and simple by
question 19 --- Gauss quadrature is the spectral theory of a
tridiagonal matrix in disguise, the finite-dimensional
shadow of \cref{ch:b3:spectral}.

\textbf{22.} Dictionary:
\begin{center}
\begin{tabular}{l|c|c|c}
 & Legendre & Hermite & Chebyshev\\
\hline
interval & $\intcc{-1}1$ & $\R$ & $\intcc{-1}1$\\
weight & $1$ & $\eu^{-t^2}$ & $(1-t^2)^{-1/2}$\\
formula & Rodrigues & $(-1)^n\eu^{t^2}
\frac{\dd^n}{\dd t^n}\eu^{-t^2}$ & $\cos(n\arccos t)$\\
norm$^2$ & $\frac2{2n+1}$ & $2^nn!\sqrt\pi$ & $\pi,
\frac\pi2$\\
habitat & quadrature & Gaussian calculus & minimax\\
\end{tabular}
\end{center}
(each with its three-term recurrence: general form for
Legendre, $H_{n+1} = 2tH_n - 2nH_{n-1}$, $T_{n+1} = 2tT_n -
T_{n-1}$). The general theory supplied what no single family
shows: reality and interlacing of roots, positivity of
quadrature weights, the mere existence of the recurrence and
of Christoffel--Darboux --- consequences of orthogonality
alone, uniform in the weight.

\textbf{23.} Existence: the linear map $\R_{2n-1}[t] \to
\R^{2n}$, $P \mapsto (P(t_1), P'(t_1), \dots, P(t_n),
P'(t_n))$, is injective (a $P$ in the kernel has $n$ double
roots and degree $\leq 2n - 1$, so $P = 0$) between spaces of
equal dimension $2n$: bijective. Pointwise error: fix $t$
not a node and choose $K$ so that $g(s) = f(s) - Hf(s) -
K\,p_n(s)^2$ vanishes at $s = t$. Then $g$ vanishes at the
$n + 1$ distinct points $t, t_1, \dots, t_n$, and $g'$
vanishes at each $t_i$ too (both $f - Hf$ and $p_n^2$ have
double zeros there). Rolle gives $n$ zeros of $g'$ strictly
between consecutive zeros of $g$ --- distinct from the nodes
--- so $g'$ has $2n$ distinct zeros; applying Rolle $2n - 1$
more times produces $\xi_t$ with $g^{(2n)}(\xi_t) = 0$.
Since $\deg Hf \leq 2n - 1$ and $p_n^2$ is monic of degree
$2n$, $g^{(2n)} = f^{(2n)} - K\,(2n)!$, whence $K =
f^{(2n)}(\xi_t)/(2n)!$ --- and the identity is trivial at
the nodes. Integration: $Q_n(f) = Q_n(Hf)$ ($Hf$ matches $f$
at the nodes) and $Q_n(Hf) = \int Hf\,w$ by exactness up to
degree $2n - 1$ (question 9), so the quadrature error is
$\int(f - Hf)\,w$. With $m, M$ the extrema of $f^{(2n)}$ on
$I$, the pointwise identity squeezes
\[
\frac{m\,h_n}{(2n)!} \;\leq\; \int_I(f - Hf)\,w
\;\leq\; \frac{M\,h_n}{(2n)!} ,
\]
and the intermediate value theorem applied to the \omterm{def:b3:topology:continuity}{continuous}
$f^{(2n)}$ delivers $\xi$. (For Legendre with $n = 2$: $h_2
= \int_{-1}^1(t^2 - \frac13)^2\dd t = \frac8{45}$, so the
error is $f^{(4)}(\xi)/135$.)

\textbf{24.} The kernel reproduces $\R_{n-1}[t]$: expanding
$q = \sum_k\frac{\langle p_k, q\rangle}{h_k}p_k$ gives
$\int_I K_n(t_i, t)\,q(t)\,w(t)\dd t = q(t_i)$ for every
$q$ of degree $\leq n - 1$. Take $q = \ell_i$: the left side
equals $\ell_i(t_i) = 1$. But $t \mapsto K_n(t_i,
t)\,\ell_i(t)$ is a polynomial of degree $\leq (n - 1) + (n
- 1) = 2n - 2$, on which $Q_n$ is exact (question 9), and it
vanishes at every node $t_j \neq t_i$ (factor $\ell_i$), so
\[
1 = \int_I K_n(t_i, t)\,\ell_i(t)\,w(t)\dd t
= w_i\,K_n(t_i, t_i)
= w_i\sum_{k=0}^{n-1}\frac{p_k(t_i)^2}{h_k} .
\]
The sum is $> 0$ (its $k = 0$ term is $1/h_0 > 0$): the
stated formula, and positivity again. Check ($n = 2$,
Legendre): $p_0 = 1$, $h_0 = 2$, $p_1 = t$, $h_1 = \frac23$;
at $t_i = \pm\frac1{\sqrt3}$,
\[
K_2(t_i, t_i) = \frac12 + \frac{1/3}{2/3} = 1,
\qquad w_i = 1,
\]
as found in question 11.

\textbf{25.} Substituting $t = \cos\theta$, the integral is
$\int_0^\pi\cos^6\theta\,\dd\theta =
\pi\,\frac{5\cdot3\cdot1}{6\cdot4\cdot2} = \frac{5\pi}{16}$
(Wallis, \cref{exo:b3:product:8}). The $n = 3$
Chebyshev--Gauss rule (question 17) has nodes
$\cos\frac\pi6 = \frac{\sqrt3}2$, $\cos\frac\pi2 = 0$,
$\cos\frac{5\pi}6 = -\frac{\sqrt3}2$ and equal weights
$\frac\pi3$:
\[
Q_3(t^6) = \frac\pi3\Bigl(2\cdot\Bigl(\frac{\sqrt3}2
\Bigr)^{6}\Bigr) = \frac\pi3\cdot\frac{54}{64}
= \frac{9\pi}{32},
\qquad
\frac{5\pi}{16} - \frac{9\pi}{32} = \frac\pi{32} .
\]
Prediction: the monic degree-$3$ orthogonal polynomial is
$2^{-2}T_3 = t^3 - \frac34t$, with $h_3 =
\frac1{16}\norm{T_3}_w^2 = \frac1{16}\cdot\frac\pi2 =
\frac\pi{32}$; and $f = t^6$ has constant $f^{(6)} = 720 =
6!$, so question 23 gives error $\frac{6!}{6!}\,h_3 =
\frac\pi{32}$ --- with no dependence on $\xi$ left, the
formula is forced to be exact, and it is.
\end{solution}
